\chapter{Introduction}

Artificial Intelligence (AI) has transformed medical image analysis by enabling
automated detection of diseases from radiological scans such as chest X-rays.
Deep learning models can support clinicians in identifying pneumonia and
related abnormalities with high speed and consistency. However, building robust
medical AI systems requires large and diverse datasets collected across
multiple hospitals. In practice, hospitals cannot freely pool patient images
into a single central repository because of legal, ethical, and security
constraints. Privacy regulations such as the Health Insurance Portability and
Accountability Act (HIPAA) and the General Data Protection Regulation (GDPR)
restrict the transfer of identifiable medical data. As a result, conventional
centralized machine learning pipelines face a fundamental conflict: high model
performance requires multi-institutional data, while privacy law requires data
localization. Federated Learning (FL) was introduced to resolve this conflict
by allowing institutions to train a shared model collaboratively while keeping
raw patient data inside each hospital~\cite{mcmahan2017communication,
guan2024federated,li2020fedprox}.

Despite strong theoretical promise, real-world medical FL deployments remain
difficult. Recent multi-center studies report that hospital firewalls,
bandwidth limits, non-Independent and Identically Distributed (non-IID) data
drift, and residual gradient leakage frequently block practical
translation~\cite{bujotzek2024real,schoenpflug2025navigating,
koutsoubis2024future,hindawi2025federated}. In resource-constrained settings
such as rural Bangladesh, these barriers are more severe because clinics often
lack high-end computing infrastructure and stable high-bandwidth networks.
Therefore, medical FL systems must be simultaneously accurate, private,
communication-efficient, and edge-deployable.

This capstone project presents FEDMED-X, a unified federated learning pipeline
for medical image classification that jointly integrates FedProx
optimization~\cite{li2020fedprox,yuan2022convergence}, Differential Privacy
via DP-SGD~\cite{abadi2016deep,mironov2017renyi,ziller2021differentially},
TopK+INT8 compression~\cite{habib2025cflsparsemed,zhou2022communication}, and
a lightweight MedNet-CBAM neural
backbone~\cite{woo2018cbam,yang2023medmnist} for edge-oriented deployment.
Rather than proposing an isolated algorithm, FEDMED-X contributes a
systems-level evaluation of the privacy--utility--communication trade-off
under realistic non-IID hospital partitions.

% -------------------------------------------------
\section{Problem Statement}
The central problem addressed in this research is the absence of a practical,
end-to-end federated medical imaging pipeline that can operate under joint
real-world constraints. Existing solutions typically optimize only one
dimension of the problem. Some methods improve non-IID
robustness~\cite{li2020fedprox,alekseenko2024distance,maleki2025handling},
some add privacy noise~\cite{abadi2016deep,ziller2021differentially}, some
compress communication payloads~\cite{habib2025cflsparsemed,asad2024fesec},
and some evaluate edge inference~\cite{vincent2026federated,
song2025edgecloud}. Very few systems evaluate all of these requirements
together under a common experimental protocol.

In clinical environments, this gap creates four concrete failures:
\begin{enumerate}[leftmargin=*]
  \item \textbf{Privacy failure:} Centralized training requires raw image
        aggregation and increases exposure to data breaches.
  \item \textbf{Heterogeneity failure:} Hospital datasets differ by patient
        population, scanner type, and disease prevalence, causing standard
        FedAvg models to drift~\cite{riyadi2026comparative}.
  \item \textbf{Communication failure:} Full-precision gradient exchange is
        expensive and often blocked by institutional
        firewalls~\cite{schoenpflug2025navigating}.
  \item \textbf{Deployment failure:} Heavy models and privacy mechanisms are
        difficult to run on low-cost edge devices suitable for rural
        clinics~\cite{shabarinath2025efficacy,memon2024edge}.
\end{enumerate}

Therefore, the research problem is not only to classify medical images, but to
design a federated system that remains usable when privacy, non-IID drift,
bandwidth limits, and edge hardware constraints are enforced at the same
time.

% -------------------------------------------------
\section{Problem Background}
Medical imaging AI has achieved strong benchmark performance under centralized
training. However, healthcare data are naturally fragmented. Each hospital
stores its own patient records, and cross-hospital transfer is restricted by
regulation and institutional policy~\cite{guan2024federated}. FL addresses
this by exchanging model updates instead of raw
images~\cite{mcmahan2017communication}. Unfortunately, recent clinical
evidence shows that FL prototypes often fail during translation. Bujotzek
\emph{et al.} demonstrated that infrastructure, governance, and
data-management bottlenecks remain major obstacles in multi-hospital
radiology FL~\cite{bujotzek2024real}. Schoenpflug \emph{et al.} reported
that network restrictions and long communication delays significantly affect
multi-country pathology FL deployments~\cite{schoenpflug2025navigating}.
Koutsoubis \emph{et al.} further emphasized that keeping data local is not
equivalent to full privacy, because shared gradients can still leak sensitive
information~\cite{koutsoubis2024future,wei2025medical}.

From an algorithmic perspective, non-IID data remains a dominant challenge.
Riyadi \emph{et al.} showed that FedProx is more robust than FedAvg under
Dirichlet label skew on MedMNIST tasks~\cite{riyadi2026comparative}. At the
same time, surveys highlight weak clinical and edge deployment benchmarks in
current medical FL literature~\cite{ghosh2026advancements,
zhou2025federated}. These findings collectively motivate a system that
combines non-IID-aware optimization, formal privacy, communication
compression, and edge-oriented model design.

% -------------------------------------------------
\section{Research Objectives}
The objectives of this study are as follows:
\begin{enumerate}[label=(\roman*),leftmargin=*]
  \item To design a federated learning framework that enables multiple
        simulated hospitals to collaboratively train a shared medical image
        classifier without exchanging raw patient data.
  \item To mitigate non-IID client drift using FedProx regularization with
        proximal coefficient $\mu = 0.01$.
  \item To enforce formal privacy protection through DP-SGD with target
        privacy budget $\epsilon = 1.0$ and $\delta = 1 \times 10^{-5}$.
  \item To reduce communication overhead by applying TopK 50\% gradient
        sparsification and INT8 quantization, targeting approximately 90\%
        bandwidth reduction.
  \item To evaluate the complete pipeline through a factorial ablation study
        and assess edge-oriented deployability using a lightweight
        MedNet-CBAM model suitable for NVIDIA Jetson Nano class hardware.
\end{enumerate}

% -------------------------------------------------
\section{Motivations}
This research is motivated by three complementary needs.

\textbf{Clinical motivation.} Bangladesh and many low-resource regions face
shortages of specialist radiologists. Affordable AI tools deployed at the edge
can support faster triage of chest X-rays in rural clinics, provided patient
privacy is preserved~\cite{vincent2026federated,shabarinath2025efficacy}.

\textbf{Technical motivation.} Most published medical FL studies evaluate
components in isolation. A practical system must quantify what happens when
privacy noise, compression, and non-IID optimization are applied
together~\cite{yesmin2025medhe,pachlasiya2026fedhealth}. Without this joint
evaluation, reported accuracies can be misleading for real deployment.

\textbf{Academic and societal motivation.} Trustworthy AI, privacy-preserving
learning, and edge intelligence are priority themes in modern computing
research and graduate study pathways. FEDMED-X is designed as a reproducible
systems benchmark that connects algorithm design with deployment reality.

% -------------------------------------------------
\section{Significance of the Research}
The significance of FEDMED-X can be summarized in four dimensions:
\begin{itemize}[leftmargin=*]
  \item \textbf{Scientific significance:} It provides a controlled
        multi-objective evaluation of accuracy, privacy budget, communication
        cost, and edge suitability under non-IID medical FL settings.
  \item \textbf{Engineering significance:} It converts known components
        (FedProx, DP-SGD, TopK, CBAM) into a coherent pipeline with explicit
        trade-off measurement.
  \item \textbf{Healthcare significance:} It supports privacy-preserving
        collaborative learning that can benefit hospitals unwilling or unable
        to share raw data.
  \item \textbf{Educational significance:} It demonstrates complex engineering
        problem-solving across machine learning, distributed systems, privacy
        mathematics, and embedded deployment constraints.
\end{itemize}

% -------------------------------------------------
\section{Research Contributions}
The major contributions of this capstone research are:
\begin{enumerate}[leftmargin=*]
  \item \textbf{Unified systems pipeline:} Development of FEDMED-X,
        integrating FedProx, DP-SGD, TopK+INT8 compression, and MedNet-CBAM
        in one reproducible federated medical classification framework.
  \item \textbf{Factorial ablation protocol:} Design and execution of a
        five-run ablation matrix (FedAvg, FedProx, FedProx+DP,
        FedProx+TopK/INT8, and full FEDMED-X) to isolate privacy and
        compression effects.
  \item \textbf{Empirical trade-off evidence:} Demonstration that FedProx
        improves non-IID robustness over FedAvg (85.58\% vs 84.46\%
        accuracy), while full privacy and compression constraints yield
        75.16\% accuracy with 90\% bandwidth reduction under
        $\epsilon = 1.0$.
  \item \textbf{Edge-oriented design:} Construction of a lightweight,
        Opacus-compatible GroupNorm-based MedNet-CBAM backbone targeting
        resource-constrained deployment scenarios such as NVIDIA Jetson
        Nano~\cite{shabarinath2025efficacy}.
  \item \textbf{Deployment-aware framing:} Positioning the work as a systems
        evaluation contribution rather than an unsupported absolute ``first
        invention'' claim, aligned with current literature gaps in edge and
        clinical translation~\cite{ghosh2026advancements,zhou2025federated}.
\end{enumerate}

% -------------------------------------------------
\section{Complex Engineering Analysis}
FEDMED-X qualifies as a complex engineering project because it cannot be
solved by a single textbook method. It requires concurrent reasoning over
conflicting objectives, multi-disciplinary knowledge, and interdependent
subsystems.

\subsection{Complex Engineering Problems}
The project maps to the following complex engineering problem attributes:
\begin{itemize}[leftmargin=*]
  \item \textbf{Depth of knowledge:} The solution requires machine learning,
        distributed optimization, differential privacy, medical data
        handling, and edge computing.
  \item \textbf{Conflicting requirements:} Higher privacy noise and stronger
        compression generally reduce accuracy; higher accuracy often
        increases communication and compute cost.
  \item \textbf{Depth of analysis:} Multiple design alternatives exist
        (FedAvg vs FedProx, DP vs no DP, compressed vs uncompressed
        updates), requiring controlled ablation rather than one-shot
        implementation.
  \item \textbf{Unfamiliarity:} Formal DP accounting~\cite{mironov2017renyi},
        non-IID federated optimization~\cite{yuan2022convergence}, and edge
        deployment constraints are beyond routine undergraduate coursework.
  \item \textbf{Stakeholder diversity:} Patients, clinicians, hospital IT
        units, regulators, and rural health workers impose different success
        criteria.
  \item \textbf{Interdependence:} Model architecture, privacy engine,
        compression module, and communication protocol affect one another
        and must be calibrated jointly.
\end{itemize}

\subsection{Complex Engineering Activities}
The complex engineering activities performed in this project include:
\begin{itemize}[leftmargin=*]
  \item \textbf{Resource integration:} Combining Colab GPU training,
        MedMNIST datasets, PyTorch/Opacus tooling, and edge-oriented model
        constraints.
  \item \textbf{Trade-off resolution:} Measuring privacy tax, compression
        effects, and joint utility loss through factorial experiments.
  \item \textbf{Innovation through synthesis:} Building an end-to-end
        evaluation framework that connects algorithm choice to deployment
        metrics.
  \item \textbf{Consequence analysis:} Assessing societal impact for
        privacy-preserving healthcare AI in low-resource environments.
  \item \textbf{Extended practice:} Applying research literature synthesis,
        reproducible experimental design, and professional technical
        documentation.
\end{itemize}

% -------------------------------------------------
\section{Thesis/Project Organization}
The remainder of this report is organized as follows:
\begin{itemize}[leftmargin=*]
  \item \textbf{Chapter 2 -- Background Study} reviews related literature on
        medical FL, privacy, communication efficiency, and edge deployment,
        and analyzes the research gap.
  \item \textbf{Chapter 3 -- Research Methodology} presents the FEDMED-X
        framework, dataset setup, algorithms, and module-level design.
  \item \textbf{Chapter 4 -- Implementation and Result Analysis} describes
        the experimental environment, preprocessing, evaluation metrics,
        ablation results, and discussion.
  \item \textbf{Chapter 5 -- Standards, Constraints, and Milestones}
        discusses ethics, societal impact, sustainability, challenges,
        constraints, and project timeline.
  \item \textbf{Chapter 6 -- Conclusion} summarizes findings, limitations,
        and future research directions.
\end{itemize}

\chaptersummary{This chapter introduced the motivation and problem foundation
of FEDMED-X. Medical AI systems need multi-hospital data, but privacy laws
restrict centralized aggregation. Federated learning offers a promising
alternative, yet practical deployments still face non-IID drift, communication
barriers, residual leakage risks, and edge-resource limitations. To address
these challenges, this project develops a unified FedProx-based pipeline with
DP-SGD privacy, TopK+INT8 compression, and a lightweight MedNet-CBAM
backbone. The chapter also defined objectives, contributions, and complex
engineering justification, and outlined the structure of the remaining
report.}